\chapter{Teclado PS/2}
\label{cap:keyboard}

\section{Vis\~ao geral}

O driver \code{src/keyboard.rs} implementa o subconjunto necess\'ario para
ler teclas de um teclado PS/2 conectado ao controlador 8042. Ele combina
conceitos de dois subsistemas do Linux:

\begin{itemize}
  \item \code{drivers/input/serio/i8042.c}: acesso ao controlador e ao
        protocolo de comandos;
  \item \code{drivers/input/keyboard/atkbd.c}: interpreta\c{c}\~ao de
        scancodes e teclas modificadoras.
\end{itemize}

N\~ao h\'a IRQ1, fila de eventos, \code{input\_device}, \emph{keycode} Linux
nem \code{evdev}. O driver faz \emph{polling} do buffer de sa\'ida do 8042 e
devolve bytes ASCII diretamente ao shell.

\section{Hardware: controlador 8042}

O 8042 exp\~oe duas portas:

\begin{table}[H]
\centering
\caption{Portas do controlador 8042.}
\label{tab:ports8042}
\begin{tabular}{@{}lll@{}}
\toprule
\textbf{Porta} & \textbf{Nome} & \textbf{Uso} \\
\midrule
\code{0x60} & Dados & Leitura de scancodes; escrita de dados ao teclado. \\
\code{0x64} & Status/Comando & Leitura de status; escrita de comandos ao 8042. \\
\bottomrule
\end{tabular}
\end{table}

O registrador de status (\code{0x64}) tem os bits usados pelo driver:

\begin{table}[H]
\centering
\caption{Bits de status do 8042.}
\label{tab:status8042}
\begin{tabular}{@{}lll@{}}
\toprule
\textbf{Bit} & \textbf{M\'ascara} & \textbf{Significado} \\
\midrule
0 & \code{0x01} & Buffer de sa\'ida cheio (h\'a dado para ler) \\
1 & \code{0x02} & Buffer de entrada cheio (n\~ao enviar comando) \\
5 & \code{0x20} & Dado veio do dispositivo auxiliar (mouse) \\
\bottomrule
\end{tabular}
\end{table}

\section{Inicializa\c{c}\~ao}

A inicializa\c{c}\~ao \'e mais do que limpar o buffer. A BIOS ou o GRUB podem
deixar o controlador em um estado em que o teclado envia scancodes do
conjunto 2, enquanto o driver espera o conjunto 1, ou com a varredura
desabilitada. O procedimento implementado \'e:

\begin{table}[H]
\centering
\caption{Sequ\^encia de inicializa\c{c}\~ao do 8042.}
\label{tab:kbdinit}
\begin{tabular}{@{}clp{7.3cm}@{}}
\toprule
\textbf{\#} & \textbf{Comando} & \textbf{Efeito} \\
\midrule
1 & leituras de \code{0x60} & Descarta bytes pendentes no buffer de sa\'ida. \\
2 & \code{0xAD} em \code{0x64} & Desabilita a porta do teclado. \\
3 & \code{0x20} em \code{0x64} & Solicita leitura do \emph{configuration byte}. \\
4 & leitura de \code{0x60} & Obt\'em o byte de configura\c{c}\~ao atual
(ou usa \code{0x45} se n\~ao houver resposta). \\
5 & ajuste & Liga tradu\c{c}\~ao (bit 6), desliga IRQ1 (bit 0) e liga o clock
(limpa bit 4). \\
6 & \code{0x60} em \code{0x64} & Prepara a escrita do novo byte de
configura\c{c}\~ao. \\
7 & byte em \code{0x60} & Envia o novo \emph{configuration byte}. \\
8 & \code{0xAE} em \code{0x64} & Reabilita a porta do teclado. \\
9 & \code{0xF4} em \code{0x60} & Envia ao teclado o comando \emph{enable
scanning}. \\
10 & leitura de \code{0x60} & Consome o ACK (\code{0xFA}), se presente. \\
11 & limpeza & Descarta qualquer res\'iduo do buffer. \\
\bottomrule
\end{tabular}
\end{table}

\begin{nota}[Por que traduzir?]
O teclado PS/2 moderno costuma falar o \emph{scan code set 2}, mas o
controlador 8042 pode apresentar ao sistema o \emph{set 1} quando o bit de
tradu\c{c}\~ao est\'a ligado. O Linux tamb\'em configura essa tradu\c{c}\~ao
no i8042. Sem ela, as teclas chegam em c\'odigos diferentes dos esperados
pela tabela do driver.
\end{nota}

As esperas usam at\'e 100.000 itera\c{c}\~oes, tanto para o buffer de entrada
esvaziar quanto para o buffer de sa\'ida encher. Se o controlador n\~ao
responder, a inicializa\c{c}\~ao continua em vez de travar.

\section{Scancodes e tradu\c{c}\~ao para ASCII}

O driver trabalha com o conjunto 1. A tabela \code{KEYMAP} tem 128 posi\c{c}\~oes
e cobre letras, d\'igitos, espa\c{c}o e alguns s\'imbolos. Posi\c{c}\~oes n\~ao
mapeadas valem zero e s\~ao ignoradas.

\begin{table}[H]
\centering
\caption{Exemplos de scancodes do conjunto 1.}
\label{tab:scancodes}
\begin{tabular}{@{}llll@{}}
\toprule
\textbf{Tecla} & \textbf{Make} & \textbf{Break} & \textbf{Resultado} \\
\midrule
\code{A} & \code{0x1E} & \code{0x9E} & \code{a} / \code{A} com Shift \\
\code{1} & \code{0x02} & \code{0x82} & \code{1} / \code{!} com Shift \\
\code{Enter} & \code{0x1C} & \code{0x9C} & \code{\textbackslash n} \\
\code{Backspace} & \code{0x0E} & \code{0x8E} & \code{0x08} \\
\code{Tab} & \code{0x0F} & \code{0x8F} & \code{\textbackslash t} \\
\code{Espa\c{c}o} & \code{0x39} & \code{0xB9} & espa\c{c}o \\
\code{Shift esquerdo} & \code{0x2A} & \code{0xAA} & ativa/desativa Shift \\
\code{Shift direito} & \code{0x36} & \code{0xB6} & ativa/desativa Shift \\
\code{Tecla estendida} & \code{0xE0} & --- & prefixo; c\'odigo seguinte ignorado \\
\bottomrule
\end{tabular}
\end{table}

\subsection{Shift}

O estado de Shift \'e mantido em \code{static SHIFT: AtomicBool}. As teclas
\code{0x2A} e \code{0x36} ligam o estado; os \emph{break codes} \code{0xAA}
e \code{0xB6} desligam. A fun\c{c}\~ao \code{apply\_shift} transforma letras
em mai\'usculas e aplica o s\'imbolo secund\'ario de d\'igitos e pontua\c{c}\~ao:

\begin{lstlisting}[caption={Mapeamento com Shift.},label={lst:shift}]
b'a'..=b'z' => byte - 32,
b'1'..=b'9' => b"!@#$%^&*("[(byte - b'1') as usize],
b'0' => b')',
b'-' => b'_',
b'=' => b'+',
/* ... */
\end{lstlisting}

\subsection{Teclas estendidas}

O prefixo \code{0xE0} liga a flag \code{EXTENDED}. O pr\'oximo scancode \'e
descartado e a flag \'e limpa. Isso evita que setas, Insert, Delete, Home,
End e teclas do teclado num\'erico sejam interpretadas como letras. A
consequ\^encia \'e que essas teclas n\~ao t\^em fun\c{c}\~ao no shell.

\subsection{Mouse}

Dados do mouse compartilham o mesmo 8042. Se o bit 5 do status estiver
ligado, o byte lido \'e descartado. Isso evita que movimentos do mouse
apare\c{c}am como letras.

\section{Fluxo de leitura}

\begin{figure}[H]
\centering
\begin{tikzpicture}[
  kstep/.style={draw=rinuxblue, fill=rinuxlight, rounded corners=2pt,
               minimum width=5.1cm, minimum height=0.8cm,
               align=center, font=\scriptsize},
  dec/.style={draw=orange!70!black, fill=orange!8, diamond, aspect=2.2,
              inner sep=1pt, align=center, font=\scriptsize},
  arrow/.style={-{Latex[length=2mm]}, thick, rinuxblue},
  node distance=0.38cm
]
\node[kstep] (status) {L\^e status (\code{0x64})};
\node[dec, below=of status] (full) {bit 0 setado?};
\node[kstep, below=of full] (read) {L\^e dado (\code{0x60})};
\node[dec, below=of read] (mouse) {bit 5 (mouse)?};
\node[dec, below=of mouse] (ext) {\code{0xE0} ou extended?};
\node[dec, below=of ext] (break) {bit 7 (break)?};
\node[kstep, below=of break] (key) {Traduz e retorna byte};
\node[kstep, right=2.4cm of full] (none) {Retorna \code{None}};
\draw[arrow] (status) -- (full);
\draw[arrow] (full) -- node[right,font=\tiny]{sim} (read);
\draw[arrow] (full.east) -- node[above,font=\tiny]{n\~ao} (none.west);
\draw[arrow] (read) -- (mouse);
\draw[arrow] (mouse) -- node[right,font=\tiny]{n\~ao} (ext);
\draw[arrow] (ext) -- node[right,font=\tiny]{n\~ao} (break);
\draw[arrow] (break) -- node[right,font=\tiny]{n\~ao} (key);
\draw[arrow] (mouse.east) -- ++(0.9,0) |- (none.east);
\draw[arrow] (ext.east) -- ++(0.5,0) |- (none.east);
\draw[arrow] (break.east) -- ++(0.2,0) |- (none.east);
\end{tikzpicture}
\caption{Fluxo de \code{read\_char()}. Caminhos de descarte levam a \code{None}.}
\label{fig:kbdflow}
\end{figure}

O shell chama \code{keyboard::read\_char()} repetidamente. Quando a fun\c{c}\~ao
retorna \code{None}, o shell tenta a serial e, se tamb\'em n\~ao houver dado,
executa \code{spin\_loop()}.

\section{API e fun\c{c}\~oes internas}

\begin{table}[H]
\centering
\caption{Fun\c{c}\~oes de \code{keyboard.rs}.}
\label{tab:kbdapi}
\begin{tabularx}{\linewidth}{@{}lX@{}}
\toprule
\textbf{Fun\c{c}\~ao} & \textbf{Descri\c{c}\~ao} \\
\midrule
\code{init()} & Configura tradu\c{c}\~ao, IRQ, clock e varredura do 8042. \\
\code{read\_char()} & L\^e e traduz um scancode; retorna \code{Option<u8>}. \\
\code{build\_keymap()} & Constr\'oi a tabela de 128 entradas em tempo de
compila\c{c}\~ao. \\
\code{apply\_shift(byte)} & Aplica o efeito de Shift a um caractere. \\
\code{wait\_input\_clear()} & Espera o buffer de entrada do 8042 esvaziar. \\
\code{wait\_output\_full()} & Espera um byte ficar dispon\'ivel na sa\'ida. \\
\code{flush\_output()} & Descarta at\'e 64 bytes pendentes. \\
\code{send\_command(cmd)} & Envia um comando ao 8042 (porta \code{0x64}). \\
\code{send\_data(data)} & Envia um byte ao teclado (porta \code{0x60}). \\
\bottomrule
\end{tabularx}
\end{table}

\section{Limita\c{c}\~oes}

\begin{itemize}
  \item \textbf{Polling}: n\~ao h\'a IRQ1; o shell precisa consultar
        constantemente. A CPU fica ocupada mesmo sem entrada.
  \item \textbf{Sem Caps Lock, Num Lock e Scroll Lock}: os LEDs e os estados
        n\~ao s\~ao tratados.
  \item \textbf{Sem AltGr}: layouts que dependem de AltGr n\~ao produzem
        caracteres adicionais.
  \item \textbf{Sem teclas estendidas funcionais}: setas, Delete, Home, End,
        Page Up/Down e teclado num\'erico s\~ao descartados.
  \item \textbf{Sem layout ABNT2}: a tabela \'e baseada em US/QWERTY; a
        posi\c{c}\~ao f\'isica das teclas e o conjunto de s\'imbolos podem
        diferir.
  \item \textbf{Sem \emph{key release} \'util}: al\'em de Shift, os
        \emph{break codes} s\~ao ignorados; n\~ao h\'a detec\c{c}\~ao de tecla
        pressionada para jogos ou atalhos.
  \item \textbf{Sem tratamento de erros}: bytes de erro (\code{0x00},
        \code{0xFF}, \code{0xEE}) n\~ao t\^em tratamento espec\'ifico.
  \item \textbf{Depend\^encia de hardware}: assume 8042 compat\'ivel. Em
        sistemas apenas USB-HID, o teclado pode n\~ao funcionar sem
        \emph{legacy USB} habilitado na BIOS.
\end{itemize}
